\section{Execution Infrastructure and Workflow}
\label{sec:execution-workflow}

The framework is meant to be used repeatedly during agent development: run suites locally or in CI, inspect failures, change the agent, and compare runs over time. Persistence and the web UI support that loop; the CLI remains the quickest path for a single failing scenario.

\subsection{Command-Line Interface}

The entry point \texttt{agent-spec-kit run PATH} discovers scenarios, resolves fixtures, executes jobs (optionally in parallel with \texttt{-n}), and writes results under \texttt{.agent\_spec\_kit/}. Runs can be grouped with \texttt{--experiment NAME} so later comparisons treat them as comparable buckets. Commands \texttt{runs} and \texttt{show RUN\_ID} list stored runs and summarise failures without opening the browser.

When an assertion fails, the terminal prints a counterexample panel aligned with the web failure card: headline, location detail, path, expected and actual summaries, and a Rich event tree. Figure~\ref{fig:cli-failure-panel} is a placeholder for a captured CLI failure panel.

\begin{figure}[htbp]
  \centering
  \fbox{\parbox{0.92\textwidth}{\centering\vspace{2.5cm}
  \textbf{Placeholder: CLI failure panel}\\
  \small \texttt{agent-spec-kit run PATH}
  \vspace{2.5cm}}}
  \caption{CLI counterexample output (placeholder): same core fields as the web FailureCard.}
  \label{fig:cli-failure-panel}
\end{figure}

Each run records git metadata, command line, and gzipped transcript and counterexample blobs in SQLite. That makes failures reviewable long after the terminal session ends.

\subsection{Web Results Browser}

The command \texttt{agent-spec-kit ui} serves a read-only React application over the local store. Figure~\ref{fig:ui-runs-list} shows the runs list with experiment and status filters. Figure~\ref{fig:ui-run-detail} shows the scenario-and-repeat table for one run; clicking a row opens the trace drawer in Figure~\ref{fig:ui-trace-drawer} with a failure card and collapsible event tree. JSON fields on tool nodes can be toggled without leaving the page.

\begin{figure}[htbp]
  \centering
  \fbox{\parbox{0.88\textwidth}{\centering\vspace{2cm}
  \textbf{Placeholder: Runs list (\texttt{/})}
  \vspace{2cm}}}
  \caption{Runs list page (placeholder).}
  \label{fig:ui-runs-list}
\end{figure}

\begin{figure}[htbp]
  \centering
  \fbox{\parbox{0.88\textwidth}{\centering\vspace{2cm}
  \textbf{Placeholder: Run detail (\texttt{/runs/:runId})}
  \vspace{2cm}}}
  \caption{Run detail: scenario and repeat rows (placeholder).}
  \label{fig:ui-run-detail}
\end{figure}

\begin{figure}[htbp]
  \centering
  \fbox{\parbox{0.88\textwidth}{\centering\vspace{2cm}
  \textbf{Placeholder: Trace drawer with FailureCard and TraceTree}
  \vspace{2cm}}}
  \caption{Trace drawer for a failing repeat (placeholder).}
  \label{fig:ui-trace-drawer}
\end{figure}

Fuzz-heavy runs also expose a trials panel so individual generated dialogues can be opened with the same trace machinery. Figure~\ref{fig:ui-fuzz-trials} is a placeholder for that view when a run includes \texttt{fuzz\_conversation} steps.

\begin{figure}[htbp]
  \centering
  \fbox{\parbox{0.88\textwidth}{\centering\vspace{2cm}
  \textbf{Placeholder: Fuzz trials panel on run detail}
  \vspace{2cm}}}
  \caption{Fuzz trials table with per-trial trace links (placeholder).}
  \label{fig:ui-fuzz-trials}
\end{figure}

The UI is read-only over the local SQLite store started by \texttt{agent-spec-kit ui --open}. It complements Langfuse-style production tracing (Chapter~\ref{chap:related-work}) but stays assertion-first: the FailureCard mirrors the CLI counterexample fields rather than asking the developer to infer expectations from spans alone.

\subsection{Compare Page and Regression Workflow}
\label{sec:ui-compare}

After changing an agent, rerunning the same scenarios under a new experiment name and opening the Compare page answers a practical question: which cases flipped between pass and fail? The route \texttt{/compare?experiments=baseline\_v1,baseline\_v2} syncs experiment slots to the query string so a comparison view can be bookmarked or shared within a team. Rows align on scenario keys; columns show the latest repeat status per experiment; cells with differing pass/fail are highlighted so regressions stand out without opening every trace. Clicking a cell opens the trace drawer for that experiment's repeat on the scenario, reusing the same FailureCard and TraceTree components as run detail.

\begin{figure}[htbp]
  \centering
  \fbox{\parbox{0.92\textwidth}{\centering\vspace{2.8cm}
  \textbf{Placeholder: Compare page}\\
  \small Two experiments; mismatched cells highlighted
  \vspace{2.8cm}}}
  \caption{Compare view across experiments (placeholder).}
  \label{fig:ui-compare}
\end{figure}

That view closes the loop described in Section~\ref{sec:counterexamples}: diagnose with a counterexample, edit the agent, rerun, and confirm at a glance whether the fix helped or regressed other scenarios. Scenarios that exist only as generative fuzz keys may be excluded from comparison when the API reports them as non-comparable across experiments; the UI surfaces that case with a short banner rather than silently showing misleading cells.
